\documentclass[10pt,a4paper]{article}
\usepackage{amsmath,amssymb,amsthm,graphicx}
\usepackage{hyperref}
\usepackage{latexsym}
\usepackage{amssymb}
\usepackage{amsthm}
\usepackage{ulem}
\usepackage{enumitem}
\setlist{itemsep=0.4pt,topsep=0pt}
\parindent 0pt
\parindent 0pt
\parskip 12pt
\newcommand{\HCF}{\mathrm{hcf}}
\newcommand{\ZZ}{\mathbb{Z}}
\newcommand{\rk}{\mathrm{rk}}
\DeclareMathOperator{\len}{length}
\DeclareMathOperator{\bl}{b}
\DeclareMathOperator{\cl}{cl}
\DeclareMathOperator{\di}{d}
\DeclareMathOperator{\Cl}{Cl}
\DeclareMathOperator{\acl}{acl}
\DeclareMathOperator{\tp}{tp}
\DeclareMathOperator{\Th}{Th}
\newcommand{\Cor}{\mathcal{C}_\mathrm{or}}
\newcommand{\Cun}{\mathcal{C}_\mathrm{un}}
\newcommand{\mons}{\mathbb{M}}
\newcommand{\strong}{\mathrm{strong}}
\newcommand{\axioms}{\Delta_{\mathrm{FSTS}}}
\newtheoremstyle{note}% hnamei
{8pt}% hSpace abovei
{14pt}% hSpace belowi
{\normalfont}% hBody fonti
{}% hIndent amounti1
{\bfseries}% hTheorem head fonti
{ }% hPunctuation after theorem headi
{.5em}% hSpace after theorem headi2
{}% hTheorem head spec (can be left empty, meaning ‘normal’)i
\theoremstyle{note}
\newtheorem{thm}{Theorem}[section]
%\newtheorem{theorem}{Theorem}
\newtheorem{lemma}[thm]{Lemma} 
\newtheorem{df}[thm]{Definition}
\newtheorem{rmk}[thm]{Remark} 
\newtheorem{prop}[thm]{Proposition}
\newtheorem{fact}[thm]{Fact}
\newtheorem{cor}[thm]{Corollary}
\newtheorem{conj}[thm]{Conjecture}
\newtheorem{ex}[thm]{Example}
\newtheorem{qu}[thm]{Question}
\newtheorem*{notation}{Notation} 
\title{The subsystem generated by an elementary subsystem is elementary, with an explicit ordering}
\author{Silvia Barbina, Enrique Casanovas}
\date{arxiv:2411.13723v3; Creative Commons Attribution 4.0 International}
\begin{document}
\maketitle
\noindent\textit{Editorial note.} The counted statement is the author's environment at source0.tex lines 909-910 and its proof environment opens at line 911; the pairing is by adjacency and the retained proof range is the whole environment.\par\medskip
\noindent\textit{Editorial note.} The preamble is the author's own package lines, macro block and theorem-environment declarations (source lines 4-65); the author's optioned amsmath line and the fleqn-dependent layout line are not retained.\par\medskip
\noindent\textit{Editorial note.} The source carries no label-number dump and no counter values are imposed, so the printed numbering is the wrapper's default; identity is the author label, or the pinned source position where the author gave none.\par\medskip
\noindent\textit{Editorial note.} Metadata provenance: the source record of the pinned copy and the licence evidence file are both present; the licence is CC BY 4.0 with status supported, and the card's source\_version matches the value recorded in that evidence file.\par\medskip
\noindent\textit{Editorial scope.} One original result of the pinned author manuscript, counted once: the generated subsystem of an elementary subsystem is elementary (author label May\_26\_2024\_2).\par\medskip
\noindent\textit{Allowed standard prerequisites (editorial).}\begin{itemize}
\item Elementary subsystem. A is elementary in M when the two structures satisfy the same formulas with parameters in A.. It is the hypothesis and part of the conclusion.
\item Generated subsystem. The subsystem generated by A is the closure of A under the partial operation.. It is the subsystem shown to be elementary.
\item Ordering extending a given one. An HF-ordering of the generated subsystem can be chosen so that its closure of A is A itself.. It is the extra information the proof provides.
\end{itemize}\medskip
	\begin{df} \label{def5_1} For $A, B \in \Cun$,
$$A \leq B \  \Leftrightarrow \ A\subseteq B \mbox{ and there is an HF-ordering of $B$ where $A$ is an initial segment}. $$
When $A \leq B$, we say that $A$ is \textbf{strong} in $B$.
\end{df}

\begin{lemma}\label{order_closure}
	If $<$ is an HF-ordering of $M$,  $A\subseteq M$  and $\cl_<(A)=A$, then  $\cl_<(\langle A\rangle )= \langle A \rangle$.
\end{lemma}

\begin{lemma} \label{May_26_2024_2} Let $M$ be an unconfined STS.  If $A\leq M$ (as in Definition~\ref{def5_1}), then  $\langle A\rangle \leq M$ and there is an HF-ordering of $\langle A\rangle$ over $A$ whose final segment $\langle A\rangle \smallsetminus A$ is well-ordered. If $N$ is an unconfined STS and $B\leq N$ and $f:A\cong B$ is an isomorphism of partial STSs, then $f$ can be extended to an isomorphism $f^\prime :\langle A\rangle \cong \langle B \rangle$.
\end{lemma}

\begin{proof} We could use Lemma~\ref{order_closure} to verify the first assertion, but the following proof also gives the isomorphism. Decompose $\langle A\rangle\smallsetminus A$ into a disjoint union $\bigcup_{n<\omega}A_n$,  where $A_0$ is the set of elements of $\langle A\rangle\smallsetminus A$ which are a product of two elements of $A$,  and $A_{n+1}$ is the set of elements of  $\langle A\rangle\smallsetminus (A\cup A_0\cup \ldots\cup A_n)$  which are a product of two elements of  $A\cup A_0\cup\ldots\cup A_n$. Fix an HF-ordering of $M$ with initial segment $A$. There are no blocks between elements of $A_0$, and the only blocks of elements of $A_0$ with smaller elements are blocks with elements of $A$ that justify their inclusion in $A_0$. Hence we can move  $A_0$ down, just above $A$,  and we can give $A_0$ a well-founded HF-ordering over $A$. This shows that $A\cup A_0\leq M$. By a similar argument we can move down each $A_{n+1}$ above $A\cup A_0\cup \ldots \cup A_n$ and give it a well-founded HF-ordering over this set. The sequence of these well-founded orders is a well-founded HF-order of $\langle A\rangle \smallsetminus A$  which makes  $\langle A\rangle$ an initial segment of the order of $M$. If we perform a similar decomposition $\bigcup_{n<\omega}B_n$ of $\langle B\rangle\smallsetminus B$, then we can  extend $f$ inductively to an isomorphism of partial Steiner triple systems  $f_n : A\cup A_0 \cup \ldots \cup A_n\cong B\cup B_0\cup\ldots \cup B_n$. The union of this ascending chain of isomorphisms is the required isomorphism $f^\prime$.
\end{proof}

\end{document}
